% concatenated from main.tex
\documentclass[11pt]{amsart}
\usepackage{hyperref}
\usepackage[capitalize, nameinlink]{cleveref}
\crefname{theorem}{Theorem}{Theorems}
\crefname{proposition}{Proposition}{Propositions}
\crefname{observation}{Observation}{Observations}
\crefname{lemma}{Lemma}{Lemmas}
\crefname{claim}{Claim}{Claims}
\crefname{problem}{Problem}{Problems}
\crefname{conjecture}{Conjecture}{Conjectures}
\crefname{question}{Question}{Questions}
\crefname{example}{Example}{Examples}
\crefname{fact}{Fact}{Facts}

%% Packages

\usepackage{amssymb}
\usepackage[bbgreekl]{mathbbol}
\usepackage{graphicx}
\usepackage{ifthen}
\usepackage{pict2e}
\usepackage{xargs}
\usepackage{xspace}
\usepackage{xcolor}
\usepackage{pgf,tikz}
\usepackage{pgfplots}
\usepackage{environ}
\usepackage{caption}
\usetikzlibrary{arrows}
\usepackage{enumitem}
\usepackage{xifthen}
\usepackage{comment}
\sloppy

\newcounter{dummy}
\makeatletter
\newcommand\myitem[1][]{\item[#1]\refstepcounter{dummy}\def\@currentlabel{#1}}
\makeatother

\makeatletter
\newsavebox{\measure@tikzpicture}
\NewEnviron{scaletikzpicturetowidth}[1]{%
	\def\tikz@width{#1}%
	\def\tikzscale{1}\begin{lrbox}{\measure@tikzpicture}%
		\BODY
	\end{lrbox}%
	\pgfmathparse{#1/\wd\measure@tikzpicture}%
	\edef\tikzscale{\pgfmathresult}%
	\BODY
}
\makeatother

%% Fonts
\DeclareSymbolFontAlphabet{\mathbb}{AMSb}

%% Macros
\newcommand{\thistheoremname}{}
\newtheorem*{genericthm*}{\thistheoremname}
\newenvironment{namedthm*}[1]
{\renewcommand{\thistheoremname}{#1}%
	\begin{genericthm*}}
	{\end{genericthm*}}

% Mathematical macros
\newcommand{\absolutevalue}[1]{|#1|}
\newcommand{\ADR}{\mathtt{AD}_{\R}}
\DeclareMathOperator{\appr}{appr}
\newcommand{\Bairespace}[1][]{
	\ifthenelse{\equal{#1}{}}{\functions{\N}{\N}}{\functions{#1}{\N}}
}
\newcommand{\bbL}{\mathbb{L}}
\newcommand{\bbX}{\mathbb{X}}
\newcommand{\Borelhomomorphism}{\le_B}
\newcommand{\Cantorspace}[1][]{
	\ifthenelse{\equal{#1}{}}{\functions{\N}{2}}{\functions{#1}{2}}
}
\newcommand{\completegraph}[1]{K_{#1}}
\newcommandx{\concatenation}[2][1 = undefined, 2 = undefined]{
	\ifthenelse{\equal{#1}{undefined}}{{}\smallfrown}{
		\ifthenelse{\equal{#2}{undefined}}{\bigoplus #1}{\bigoplus_{#1} #2}
	}
}
\newcommand{\constantsequence}[2]{\sequence{#1}^{#2}}
\newcommand{\continuoushomomorphism}{\le_c}
\newcommand{\definedterm}[1]{\emph{#1}}
\newcommand{\domain}[1]{\mathrm{dom}(#1)}
\newcommand{\from}{\colon}
\newcommandx{\functions}[3][3 =]{
	\ifthenelse{\equal{#3}{}}{#2^{#1}}{#2_{#3}^{#1}}
}
\newcommand{\goesto}{\rightarrow}
\newcommand{\graphmetric}[1]{d_{#1}}
\newcommand{\Gzero}[1][]{
	\ifthenelse{\equal{#1}{}}
	{\mathbb{G}_0}
	{\mathbb{G}_{0,n}}
}
\newcommandx{\Hzero}[2][2 = undefined]{
	\ifthenelse{\equal{#2}{undefined}}
	{\mathbb{H}_{#1}}
	{\mathbb{H}_{#1, #2}}
}
\newcommandx{\intersection}[2][1 =, 2 =]{
	\ifthenelse{\equal{#1}{}}{\cap}{
		\ifthenelse{\equal{#2}{}}{\bigcap #1}{{\bigcap_{#1} #2}}
	}
}
\newcommand{\Lzero}[1][]{\ifthenelse{\equal{#1}{}}{\bbL_0}{L_{0, #1}}}
\newcommand{\Lzerospace}[1][]{\ifthenelse{\equal{#1}{}}{\bbX_0}{X_{0, #1}}}
\newcommand{\mathand}{\text{ and }}
\newcommand{\modulo}[1]{\ (\text{mod } 2)}
\newcommand{\N}{\mathbb{N}}
\newcommand{\odometer}[1]{\mathbb{\Sigma}_{#1}}
\newcommand{\pair}[2]{(#1, #2)}
\newcommand{\partialto}{\rightharpoonup}
\newcommand{\pathlength}[1]{\mathrm{length}(#1)}
\newcommandx{\product}[2][1 =, 2 =]{
	\ifthenelse{\equal{#1}{}}{\times}{
		\ifthenelse{\equal{#2}{}}{\prod #1}{{\prod_{#1} #2}}
	}
}

\newcommand{\rbar}{\bar{r}}
\newcommand{\rr}{\textcolor{red}}
\newcommandx{\sequence}[2][2 = undefined]{
	\ifthenelse{\equal{#2}{undefined}}{(#1)}{
		(#1)_{#2}
	}
}
\newcommand{\sbar}{\bar{s}}
\newcommandx{\set}[2][2 = undefined]{
	\ifthenelse{\equal{#2}{undefined}}{\{ #1 \}}{
		\{ #1 \suchthat #2 \}
	}
}
\newcommandx{\sets}[3][3 =]{
	\ifthenelse{\equal{#3}{}}{[#2]^{#1}}{[#2]^{#1}_{#3}}
}
\newcommand{\shift}[2]{\mathbb{S}_{#1, #2}}
\newcommand{\signedlength}[2]{\length{#1}_{#2}}
\newcommand{\suchthat}{\mid}
\newcommand{\R}{\mathbb{R}}
\renewcommand{\restriction}[2]{#1 \upharpoonright #2}

\newcommand{\triple}[3]{(#1, #2, #3)}
\newcommandx{\union}[2][1 =, 2 =]{
	\ifthenelse{\equal{#1}{}}{\cup}{
		\ifthenelse{\equal{#2}{}}{\bigcup #1}{{\bigcup_{#1} #2}}
	}
}
\newcommand{\undirectable}[1]{\mathbb{L}_{#1}}
\newcommand{\Xzero}[1]{X_{#1}}

% Noun macros

\newcommand{\Baire}{Baire\xspace}
\newcommand{\Borel}{Bor\-el\xspace}
\newcommand{\Hausdorff}{Haus\-dorff\xspace}
\newcommand{\Kechris}{Kech\-ris\xspace}
\newcommand{\Polish}{Po\-lish\xspace}
\newcommand{\Solecki}{So\-lec\-ki\xspace}
\newcommand{\Todorcevic}{To\-dor\-\v{c}ev\-ic\xspace}

% Picture macros
\renewcommand{\Box}{
	\begin{picture}(6,6)
	\put(0,0){\line(1,0){6}}
	\put(0,0){\line(0,1){6}}
	\put(6,0){\line(0,1){6}}
	\put(0,6){\line(1,0){6}}
	\put(0,0){\line(1,1){6}}
	\put(6,0){\line(-1,1){6}}
	\end{picture}
}

% Proof environment
\newenvironment{theoremproof}{
	\renewcommand{\qedsymbol}{$\Box$}
	\begin{proof}
	}{\end{proof}}

% Theorem environments
\newtheorem{theorem}{Theorem}[section]
\newtheorem{lemma}[theorem]{Lemma}
\newtheorem{Claim}[theorem]{Claim}
\newtheorem{claim}[theorem]{Claim}
\newtheorem{subclaim}[theorem]{Subclaim}
\newtheorem{corollary}[theorem]{Corollary}
\newtheorem{proposition}[theorem]{Proposition}
\newtheorem{conjecture}[theorem]{Conjecture}
\newtheorem{question}[theorem]{Question}
\newtheorem{problem}[theorem]{Problem}
\newtheorem{fact}[theorem]{Fact}
\newtheorem{ax}{Axiom}


\theoremstyle{definition}

\newtheorem{definition}[theorem]{Definition}
\newtheorem{example}[theorem]{Example}
\newtheorem{remark}[theorem]{Remark}

\numberwithin{equation}{section}

% More mathematical macros
%\newcommand{\R}{\mathbb{R}}
\newcommand{\Q}{\mathbb{Q}}
%\newcommand{\N}{\mathbb{N}}
\newcommand{\Z}{\mathbb{Z}}
\newcommand{\rank}{\qopname\relax o{rk}}
\newcommand{\dlength}{\qopname\relax o{dilength}}
\newcommand{\length}{\qopname\relax o{length}}
\newcommand{\iB}{\mathcal{B}}
\newcommand{\iF}{\mathcal{F}}
\newcommand{\iU}{\mathcal{U}}
\newcommand{\iV}{\mathcal{V}}
\newcommand{\al}{\alpha}
\newcommand{\la}{\lambda}
\newcommand{\om}{\mathbb{N}}
\newcommand{\omm}{[\mathbb{N}]^\mathbb{N}}
\newcommand{\oom}{\mathbb{N}^\mathbb{N}}
\newcommand{\si}{\sigma}
\newcommand{\eps}{\varepsilon}
\DeclareMathOperator{\maxdeg}{maxdeg}
\newcommand{\RR}{\mathbb{R}}
\newcommand{\bd}{\begin{definition}}
	\newcommand{\ed}{\end{definition}}
\newcommand{\solc}{\mathop{sc}}
\newcommand{\lab}{\label}
\DeclareMathOperator{\graph}{graph}
\DeclareMathOperator{\dom}{dom}
\DeclareMathOperator{\proj}{proj}
\DeclareMathOperator{\ran}{ran}
\DeclareMathOperator{\imm}{im}
\DeclareMathOperator{\rd}{rd}
\DeclareMathOperator{\supp}{supp}
\DeclareMathOperator{\mgs}{mgs}
\DeclareMathOperator{\dist}{dist}
\DeclareMathOperator{\didistance}{didist}
\newcommand{\mc}{\mathcal}
\newcommand{\mb}{\mathbf}
\newcommand{\bs}{\mathbf{\Sigma}^1_1}
\newcommand{\bbg}{\mathbf{\Gamma}}
\newcommand{\bbo}{\mathbf{\Delta}^1_1}
\newcommand{\bp}{\mathbf{\Pi}^1_1}
\newcommand{\ls}{\Sigma^1_1}
\newcommand{\lbo}{\Delta^1_1}
\newcommand{\lp}{\Pi^1_1}
\newcommand{\ve}{\varepsilon}
\newcommand{\mf}{\mathfrak}
\newcommand{\nr}{\lnot R}
\newcommand{\id}{id}
\newcommand{\distance}[3]{\ifthenelse{\isempty{#3}}{\dist(#1,#2)}{\dist^{#3}(#1,#2)}}
\newcommand{\didist}[3]{\ifthenelse{\isempty{#3}}{\didistance(#1,#2)}{\didistance^{#3}(#1,#2)}}
\newcommand{\digraph}[3]{\ifthenelse{\equal{#1}{b}}{\mathbb{#2}_{#3}}
	{{#2}_{#3}}}
\newcommand{\linegraph}[3]{\ifthenelse{\equal{#1}{b}}{\mathbb{#2}_{#3}}
	{#2_{#3}}}
\newcommand{\specialvertex}[2]{s_{#1}}
\newcommand{\underlyingspace}[3]{\ifthenelse{\equal{#1}{b}}{\mathbb{#2}_{#3}}
	{#2_{#3}}}
\newcommand{\distanceset}[2]{\ifthenelse{\isempty{#2}}{D(#1)}{D^{#2}(#1)}}
\newcommand{\Borelchromatic}[1]{\chi_B(#1)}
\newcommand{\restrcitedto}[2]{#1|#2}
\newcommand{\connecteder}[1]{E_{#1}}
%\newcommand{\dist}{dist}
%\newcommand{\length}{length}
%\newcommand{\dlength}{dilength}
\newcommand{\concatt}{%
	\mathbin{\raisebox{1ex}{\scalebox{.7}{$\frown$}}}%
}
\newcommand{\ovl}{\overline}
\newcommand{\pred}{pred}
\newcommand{\directed}[1]{\mathbb{L}_{#1}}
\newcommand{\Hzerodirected}[2]{\mathbb{H}_{#1,#2}}
\newcommand{\aStage}[1]{n^{#1}}
\newcommand{\aCode}[1]{f^{#1}}
\newcommand{\aCodeB}[1]{c^{#1}}
\newcommand{\aCodeD}[1]{d^{#1}}
\newcommand{\aMapV}[1]{\phi^{#1}}
\newcommand{\aMapE}[2]{\psi^{#1}_{#2}}
\newcommand{\aBoundStage}[1]{\aCodeB{#1},\aStage{#1}}
\newcommand{\aDifferent}[2]{\aCodeB{#2},\aStage{#1},\aStage{#2}}
\newcommand{\aCodeStage}[1]{\aCode{#1},\aStage{#1}}
\newcommand{\Lias}[3]{((#2_{i})_{i<#1},(#3_{i})_{0<i<#1})}
\newcommand{\aLias}[2]{#1^{#2}}
\newcommand{\imageConf}[2]{A(#1,#2)}
\newcommand{\BorelInvariant}[2]{B(#1,#2)}
\newcommand{\twocoloring}[2]{c_{#1,#2}}
\newcommand{\oddpair}[1]{odd $#1$-pair}
\newcommand{\lgth}{\length}
%\newcommand{\N}{\mathbb{N}}
\newcommand{\Ramseyspace}{[\N]^{\N}}
%\newcommand{\dist}{dist}
\newcommand{\Shiftmap}{\mathcal{S}}
\newcommand{\Shiftgraph}{\mathcal{G}_\Shiftmap}
\newcommand{\Shiftgraphmodified}{\mathcal{G}'_{\Shiftmap}}
\newcommand{\Homomorphism}[2]{Hom(#1,#2)}
\newcommand{\Homomorphisma}[2]{Hom^*(#1,#2)}
\newcommand{\Tree}[1]{T_{#1}}
\newcommand{\Homothree}[1]{Hom(\Tree{3},#1)}
\newcommand{\Homothreea}[1]{Hom^*(\Tree{3},#1)}
\newcommand{\Deltaabsolute}[1]{\chi_{\Delta^1_2ab}(#1)}
\newcommand{\Bairemeasurablechromatic}[1]{\chi_{BM}(#1)}
\newcommand{\Measurablechromatic}[1]{\chi_{\mu}(#1)}
\newcommand{\Ztwocube}{\mathbb{Z}_2*\mathbb{Z}_2*\mathbb{Z}_2}
\newcommand{\SG}{\mathcal{G}_\mathcal{S}}



\newcommand{\rstr}{\restriction}
\newcommand{\vfi}{\varphi}
\newcommand{\sse}{\subseteq}
\newcommand{\est}{\emptyset}


\newcommand{\treegraph}[1]{Tree(#1)}
\newcommand{\fG}{\mathcal{G}}
\newcommand{\fA}{\mathcal{A}}
\newcommand{\fB}{\mathcal{B}}
\newcommand{\fC}{\mathcal{C}}
\newcommand{\fD}{\mathcal{D}}
\newcommand{\fE}{\mathcal{E}}
\newcommand{\Epsilon}{\fE}
\newcommand{\fF}{\mathcal{F}}
%\newcommand{\fG}{\mathcal{G}}
\newcommand{\fH}{\mathcal{H}}
\newcommand{\fI}{\mathcal{I}}
\newcommand{\fJ}{\mathcal{J}}
\newcommand{\fK}{\mathcal{K}}
\newcommand{\fL}{\mathcal{L}}
\newcommand{\fM}{\mathcal{M}}
\newcommand{\fN}{\mathcal{N}}
\newcommand{\fO}{\mathcal{O}}
\newcommand{\fP}{\mathcal{P}}
\newcommand{\fQ}{\mathcal{Q}}
\newcommand{\fR}{\mathcal{R}}
\newcommand{\fS}{\mathcal{S}}
\newcommand{\fT}{\mathcal{T}}
\newcommand{\fU}{\mathcal{U}}
\newcommand{\fV}{\mathcal{V}}
\newcommand{\fW}{\mathcal{W}}
\newcommand{\fX}{\mathcal{X}}
\newcommand{\fY}{\mathcal{Y}}
\newcommand{\fZ}{\mathcal{Z}}

% Tikz settings

\title{Hyperfiniteness on the set covered infinitely often by separated Borel sets of a bounded-degree decomposition}
\author{Balázs Bursics, Zoltán Vidnyánszky}
\date{arxiv:2412.01315v2; CC BY 4.0}
\begin{document}
\maketitle
\setcounter{section}{1}
\setcounter{theorem}{2}
\noindent\textit{Editorial note.} The counted statement is the introduction's Theorem t:separation2 (source lines 425-430). Its proof is not the adjacent environment: the author restates the same theorem as a namedthm* at source line 582 and proves it at 589-609, so the restatement is retained as context and makes the pairing visible; the pairing is by the author's own reference to the label t:separation2, not by proximity.\par\medskip
\noindent\textit{Editorial note.} The retained section and theorem counters restore the article's numbering, so the retained definition prints as Definition 1.3 and the counted theorem as Theorem 1.4, exactly as in the source. The Claim carried inside the proof is renumbered because the article's section headings are not reproduced.\par\medskip
\noindent\textit{Editorial note.} The footnote in the theorem statement, that graphs are not assumed irreflexive throughout the paper, is part of the author's statement and is retained verbatim.\par\medskip
\noindent\textit{Editorial note.} No retained line contains a citation, a figure or an included file, so the package declares neither a bibliography nor an asset.\par\medskip
\noindent\textit{Editorial scope.} One original result: the article's Theorem t:separation2 with its complete author proof, together with the author's own definition of a k-separated set. The article's main theorem t:ramsey hyperfin, the earlier Lemma l:indep2 and the generic-hyperfiniteness theorem of Subsection 3.3 are not counted here and are not used as premises.\par\medskip
\noindent\textit{Allowed standard prerequisites (editorial).}\begin{itemize}
\item Countable Borel equivalence relations and their graph decompositions. A CBER E on a standard Borel space X has countable classes, so E is the union of the graphs of Borel involutions, and an increasing sequence of bounded-degree Borel graphs G\_n with union E is the standard decomposition used throughout the article. Borel hyperfiniteness of a restriction E restricted to a Borel set means a Borel reduction of that restriction to eventual equality.. Both the statement and the proof start from such a decomposition; no hyperfiniteness of E itself is assumed.
\item Graph distance and k-separated sets. For a graph G and a set B of vertices, B is k-separated when any two distinct points of B have G-distance greater than k; the article's own definition is retained from source lines 419-421.. The hypothesis that each B\_n is f(n)-separated in G\_n is what bounds the diameter of the components built in the proof.
\item Elementary graph metric estimates. If a shortest path between two vertices leaves a set and the intermediate vertices lie in one connected component of a fixed subgraph, the triangle inequality bounds the distance between the endpoints by the diameter of that component plus the two entering edges.. Used inside the Claim at source line 602 to obtain dist\_\{G\_\{n+1\}\}(x,x') <= f(n)+2.
\end{itemize}\medskip
	\begin{definition}
		Let $G$ be a graph, $B \subseteq V(G)$. The set $B$ is called \emph{$k$-separated} if for any $x \neq x' \in B$ we have $\dist_{G}(x,x')>k$.
	\end{definition}

	\begin{theorem}
		\label{t:separation2}
		Let $E$ be a CBER on the space $X$, let $(G_n)_{n \in \N}$ be an increasing sequence of bounded degree Borel graphs\footnote{Throughout the paper, graphs are not necessarily assumed to be irreflexive.} such that $\bigcup_n G_n=E$, and let $f:\N\to\N$ be such that $\forall n\in\N\; f(n+1)\ge2\cdot(f(n)+1)$. Moreover, assume that $B_n \subseteq X$ are Borel sets so that every $B_n$ is $f(n)$-separated in $G_n$. 
		
		Then $\restriction{E}{B}$ is hyperfinite, where $B=\{x:\exists^\infty n \ (x \in B_n)\}$.
	\end{theorem}

	\begin{namedthm*}{Theorem \ref{t:separation2}}
		
		Let $E$ be a CBER on the space $X$, let $(G_n)_{n \in \N}$ be an increasing sequence of bounded degree Borel graphs such that $\bigcup_n G_n=E$, and let $f:\N\to\N$ be such that $\forall n\in\N\; f(n+1)\ge2\cdot(f(n)+1)$. Moreover, assume that $B_n \subseteq X$ are Borel sets so that every $B_n$ is $f(n)$-separated in $G_n$. 
		
		Then $\restriction{E}{B}$ is hyperfinite, where $B=\{x:\exists^\infty n \ (x \in B_n)\}$.
	\end{namedthm*}

	\begin{proof}
		
		We define a sequence $H_0 \subseteq H_1 \subseteq \dots \subseteq H_n \subseteq \dots \subseteq E$ inductively. Let $H_0=\emptyset$. 
		
		Given $H_n$, let \[H_{n+1}=H_{n} \cup \{(x,y) \in \restriction{G_{n+1}}{B}:\; x\in B_{n+1} \text{ or } y\in B_{n+1}\}.\]
		
		First we check that $\rstr{\bigcup_{n \in \N} H_n}{B}=\restriction{E}{B}$. Indeed, let $(x,y) \in \rstr{G_n}{B}$. Then there is some $k \geq n$ such that $x \in B_k$. But then $(x,y) \in H_k$. 
		
		\begin{claim}
			Every connected component of $H_n$ has diameter $\le f(n)$. 
			
		\end{claim}
		\begin{proof}
			We prove this by induction on $n$. The case $n=0$ is clear, so assume that we have shown the statement for $n$. We claim that if $H$ is a connected component of $H_{n+1}$ then it contains at most one point of $B_{n+1}$. Indeed, otherwise let $x \neq x' $ from $B_{n+1}$ in $V(H)$ be such that for $p=(x_0,\dots,x_k)$, the injective path with minimal length from $x_0=x$ to $x_k=x'$ in $H$ we have $x_1,\ldots,x_{k-1}\notin B_{n+1}$. Then $x_1,\ldots,x_{k-1}$ are contained in a single connected component $H'$ of $H_n$, and by the inductive hypothesis $\dist_{H_n}(x_1,x_{k-1})\le f(n).$ But this implies $\dist_{G_{n+1}}(x,x')\le f(n)+2\le f(n+1),$ a contradiction.
			
			Thus, $H$ contains at most one point, say $x$, from $B_{n+1}$. Then \[V(H)=\bigcup \{[y]_{H_{n}}:(x,y) \in G_{n+1}\}.\]
			yielding \[\text{diam}(H) \le 2\cdot (f(n)+1) \le f(n+1).\]
		\end{proof}
		
		As $H_n\sse G_n$ which has bounded degrees, this, in turn also ends the proof of Theorem \ref{t:separation2}.
	\end{proof}

\end{document}
